\documentclass[10pt,a4paper]{amsart}
\usepackage[margin=19mm]{geometry}
\usepackage{amsmath,amssymb,mathtools}
\usepackage[scr=rsfs,cal=euler]{mathalfa}
\usepackage{xfrac}
\usepackage[unicode,hidelinks,bookmarks=false]{hyperref}
\newcommand{\ie}{i.e.\ }
\newcommand{\cf}{cf.\ }
\newcommand{\resp}{respectively\ }
\newcommand{\Subsection}{\subsection}
\providecommand{\nobreakdash}{\nobreak\hbox{-}}
\setlength{\emergencystretch}{2em}
\multlinegap0pt
\usepackage{bm}
\usepackage{bbm}
\usepackage{dsfont}
\usepackage{xspace}
\usepackage{cancel}
\usepackage{mathtools}
\usepackage{amsthm}
\makeatletter
\@ifundefined{theorem}{\newtheorem{theorem}{Theorem}}{}
\@ifundefined{proposition}{\newtheorem{proposition}[theorem]{Proposition}}{}
\makeatother
\newcommand{\bg}{\bar{\mathfrak g}}
\newcommand{\mcF}{\mathcal F}
\title[A description of the depth-$r$ Bernstein center for rational]{A description of the depth-$r$ Bernstein center for rational depths}
\author{Sarbartha Bhattacharya, Tsao-Hsien Chen}
\date{Source: arXiv:2510.07845v1 (CC BY 4.0)}
\begin{document}
\maketitle

\noindent\emph{Source and licence.} A description of the depth-$r$ Bernstein center for rational depths. Version arXiv:2510.07845v1 (CC BY 4.0), distributed under the Creative Commons Attribution 4.0 International licence, as recorded in the arXiv licence metadata for this exact version and shown on the abstract page https://arxiv.org/abs/2510.07845v1. Full rights notice: licenses/arxiv-2510.07845-CC-BY-4.0.txt.\par\smallskip

\noindent\textit{Editorial scope.} Counted unit: the Proposition whose source label is \texttt{fourierprop} in the article (source0.tex lines 906--915, source label \texttt{fourierprop}), delivered together with the article's own complete original proof of it (source0.tex lines 916--953, one \texttt{proof} environment of 38 lines). No mathematics is written, repaired or re-derived: statement and proof are the article's own text line for line. Required context retained uncounted: fourierdefnlie (lines 840--842); fourierdefndual (lines 844--846). Every definition, display and label of those lines is retained unchanged. Omitted from this package and not redistributed: 1--839, 843--843, 847--905, 954--2211. Nothing inside a retained excerpt is omitted or altered: every retained body is delivered exactly as the article's lines are written, so no omission receipt of any kind accompanies this package. Citations: the retained excerpts carry no citation command at all, so no bibliography entry of the article is transcribed and no reference is redistributed. Reference adaptation: none was needed and none was made; every reference inside the delivered unit resolves inside it, so no unresolved reference or citation is produced. Printed slips kept literal: no printed word, symbol, hypothesis or number of the counted statement or of its proof was altered. Layout and class: the article's own class is \texttt{article} with packages including inputenc, amsmath, amssymb, graphicx, setspace, centernot, amsthm, fullpage, amsfonts, float, mathtools, array, multicol, enumitem; the wrapper is the lane's pinned \texttt{amsart} editorial wrapper at 10pt on A4, which restates the theorem-like environments the retained text needs and adds the article's own macro definitions that the retained text needs (\texttt{\textbackslash bg}, \texttt{\textbackslash mcF}), with extra wrapper packages (bm, bbm, dsfont, xspace, cancel, mathtools). The delivered numbering is the unit's own. Assets: no retained span contains a figure environment, a graphics-inclusion command, a PDF-page inclusion, a table or a TikZ picture; no image file is copied into this package, the package declares no asset dependency and no image file is redistributed with it. Licence: version arXiv:2510.07845v1 of this article is distributed under the Creative Commons Attribution 4.0 International licence, as recorded in the arXiv licence metadata for this exact version and shown on the abstract page https://arxiv.org/abs/2510.07845v1. A full rights notice is delivered at licenses/arxiv-2510.07845-CC-BY-4.0.txt. Characters: every retained excerpt is pure ASCII, so the delivered engine, XeLaTeX with its default Latin Modern text font, reports no missing character.
\begin{equation}\label{fourierdefnlie}
    \mcF_{\bg_{\sigma,r}^F}(f)(X^*)=\left|\bg_{\sigma,r}^F\right|^{-1/2}\sum_{Y\in \bg_{\sigma,r}^F}\Tilde\psi( X^*(Y))f(Y).
\end{equation}

\begin{equation}\label{fourierdefndual}
  \mcF_{(\bg_{\sigma,r}^F)^*}(f)(Y)=\left|(\bg_{\sigma,r}^F)^*\right|^{-1/2}\sum_{X^*\in (\bg_{\sigma,r}^F)^*}\Tilde\psi( X^*(Y))f(X^*).
\end{equation}

\begin{proposition}[Properties of Fourier transform]\label{fourierprop}
    Let  $f,g \in C\left(\bg_{\sigma,r}^F\right)$ and $\tilde f, \tilde g \in C\left((\bg_{\sigma,r}^F)^*\right)$. The Fourier transforms defined in \eqref{fourierdefnlie} and \eqref{fourierdefndual} have the following properties:
    \begin{itemize}
        \item[(i)]  $\left(\mcF_{\bg_{\sigma,r}^F}(f),\mcF_{\bg_{\sigma,r}^F}(g)\right)_{(\bg_{\sigma,r}^F)^*}= \left(f,g\right)_{\bg_{\sigma,r}^F}$. 
        \item [(ii)]  $\mcF_{\bg_{\sigma,r}^F}\circ \mcF_{(\bg_{\sigma,r}^F)^*}(\tilde f) = (\tilde f)^-$ and $ \mcF_{(\bg_{\sigma,r}^F)^*}\circ\mcF_{\bg_{\sigma,r}^F}(f)=f^-$, where $f^-(X)=f(-X)$. 
        \item[(iii)]  $\mcF_{\bg_{\sigma,r}^F}(f*g)= \mcF_{\bg_{\sigma,r}^F}(f)\cdot \mcF_{\bg_{\sigma,r}^F}(g)$ and $\mcF_{(\bg_{\sigma,r}^F)^*} (\tilde f*\tilde g )=\mcF_{(\bg_{\sigma,r}^F)^*} (\tilde f) \cdot \mcF_{(\bg_{\sigma,r}^F)^*}(\tilde g) $.
        \item[(iv)] $\mcF_{\bg_{\sigma,r}^F}(f\cdot g)= \mcF_{\bg_{\sigma,r}^F}(f)* \mcF_{\bg_{\sigma,r}^F}(g)$ and $\mcF_{(\bg_{\sigma,r}^F)^*} (\tilde f\cdot \tilde g )=\mcF_{(\bg_{\sigma,r}^F)^*} (\tilde f) * \mcF_{(\bg_{\sigma,r}^F)^*}(\tilde g)$.
    \end{itemize}
    As a consequence, we see that the Fourier transform $\mcF_{\bg_{\sigma,r}^F}:C(\bg_{\sigma,r}^F) \longrightarrow C\left((\bg_{\sigma,r}^F)^*\right)$ is an algebra isomorphism with inverse given by $\mcF_{(\bg_{\sigma,r}^F)^*}\circ \mcF_{\bg_{\sigma,r}^F}\circ \mcF_{(\bg_{\sigma,r}^F)^*}$, where multiplication on $C\left((\bg_{\sigma,r}^F)^*\right)$ is given by the usual pointwise one. The analogous result holds true for  $\mcF_{(\bg_{\sigma,r}^F)^*}:C\left((\bg_{\sigma,r}^F)^*\right) \longrightarrow C(\bg_{\sigma,r}^F) $.
\end{proposition}

\begin{proof}
    Most of these proofs follow from simple calculations using the definitions.
    \begin{itemize}
       \item[(i)]  Follows immediately from calculations.
        \item [(ii)] Let $Y \in \bg_{\sigma,r}^F$. Then, \begin{align*}
            \mcF_{(\bg_{\sigma,r}^F)^*}\circ\mcF_{\bg_{\sigma,r}^F}(f)(Y)& =\left|(\bg_{\sigma,r}^F)^*\right|^{-1/2}\sum_{X^*\in (\bg_{\sigma,r}^F)^*}\Tilde\psi( X^*(Y))\mcF_{\bg_{\sigma,r}^F}(f)(X^*)\\
            &= \left|(\bg_{\sigma,r}^F)^*\right|^{-1/2}\sum_{X^*\in (\bg_{\sigma,r}^F)^*}\Tilde\psi( X^*(Y)) \left(\left|\bg_{\sigma,r}^F\right|^{-1/2}\sum_{Y_1\in \bg_{\sigma,r}^F}\Tilde\psi( X^*(Y_1))f(Y_1)\right)\\
            &= \left|\bg_{\sigma,r}^F\right|^{-1}\sum_{X^*\in (\bg_{\sigma,r}^F)^*}\sum_{Y_1\in \bg_{\sigma,r}^F}\Tilde\psi( X^*(Y+Y_1)) f(Y_1)\\
            &= \sum_{Y_1\in \bg_{\sigma,r}^F}f(Y_1)\left( \left|\bg_{\sigma,r}^F\right|^{-1} \sum_{X^*\in (\bg_{\sigma,r}^F)^*} \Tilde\psi( X^*(Y+Y_1))\right)
        \end{align*}
        We know that the inside sum $ \frac{1}{\left|\bg_{\sigma,r}^F\right|}\sum_{X^*\in (\bg_{\sigma,r}^F)^*} \Tilde\psi( X^*(Y+Y_1))= \begin{cases} 0, & \text{ if } Y \neq -Y_1\\
        1, &\text{ if } Y = -Y_1\\
        \end{cases}$ which gives us $\mcF_{(\bg_{\sigma,r}^F)^*}\circ\mcF_{\bg_{\sigma,r}^F}(f)(Y)= f (-Y)$ and we are done. 
        \item[(iii)] Let $X^* \in (\bg_{\sigma,r}^F)^*$. Then, 
        \begin{align*}
            \mcF_{\bg_{\sigma,r}^F}(f*g)(X^*)&= \left|\bg_{\sigma,r}^F\right|^{-1/2}\sum_{Y\in \bg_{\sigma,r}^F}\Tilde\psi( X^*(Y))f*g(Y)\\
            &= \left|\bg_{\sigma,r}^F\right|^{-1/2}\sum_{Y\in \bg_{\sigma,r}^F}\Tilde\psi( X^*(Y))\left( \left|\bg_{\sigma,r}^F\right|^{-1/2}\sum_{Z\in \bg_{\sigma,r}^F}f(Y-Z)g(Z)\right)\\
            &= \left|\bg_{\sigma,r}^F\right|^{-1} \sum_{Y\in \bg_{\sigma,r}^F}\Tilde\psi( X^*(Y-Z))f(Y-Z)\sum_{Z\in \bg_{\sigma,r}^F}\Tilde\psi( X^*(Z))g(Z)\\
            &=\left(\left|\bg_{\sigma,r}^F\right|^{-1/2}\sum_{Y\in \bg_{\sigma,r}^F}\Tilde\psi( X^*(Y-Z))f(Y-Z)\right)\left(\left|\bg_{\sigma,r}^F\right|^{-1/2}\sum_{Z\in \bg_{\sigma,r}^F}\Tilde\psi( X^*(Z))g(Z)\right)\\
            &=  \mcF_{\bg_{\sigma,r}^F}(f)(X^*)\cdot  \mcF_{\bg_{\sigma,r}^F}(g)(X^*)\\
        \end{align*} 
        \item[(iv)] For any $f,g \in  C\left(\bg_{\sigma,r}^F\right)$, we see using $(i)$ that 
        \begin{equation}\label{fourierpf1}
            \mcF_{(\bg_{\sigma,r}^F)^*}\circ\mcF_{\bg_{\sigma,r}^F}(f*g)= \left(\mcF_{(\bg_{\sigma,r}^F)^*}\circ\mcF_{\bg_{\sigma,r}^F}(f)\right)*\left(\mcF_{(\bg_{\sigma,r}^F)^*}\circ\mcF_{\bg_{\sigma,r}^F}(g)\right).
        \end{equation}
        Further, applying $\mcF_{(\bg_{\sigma,r}^F)^*}\circ \mcF_{\bg_{\sigma,r}^F}\circ \mcF_{(\bg_{\sigma,r}^F)^*}$ to both sides of $(ii)$, we have 
       \[
       \mcF_{(\bg_{\sigma,r}^F)^*}\circ \mcF_{\bg_{\sigma,r}^F}\circ \mcF_{(\bg_{\sigma,r}^F)^*} \circ \mcF_{\bg_{\sigma,r}^F}(f*g)= f*g=\mcF_{(\bg_{\sigma,r}^F)^*}\circ \mcF_{\bg_{\sigma,r}^F}\circ \mcF_{(\bg_{\sigma,r}^F)^*}\left(\mcF_{\bg_{\sigma,r}^F}(f)\cdot \mcF_{\bg_{\sigma,r}^F}(g)\right)
        \]
         Then, applying $\mcF_{(\bg_{\sigma,r}^F)^*} \circ \mcF_{\bg_{\sigma,r}^F}$ to both sides, we have by using \eqref{fourierpf1}
        \begin{align*}
            \mcF_{(\bg_{\sigma,r}^F)^*} \circ \mcF_{\bg_{\sigma,r}^F}(f*g)&=\mcF_{(\bg_{\sigma,r}^F)^*}\left(\mcF_{\bg_{\sigma,r}^F}(f)\cdot \mcF_{\bg_{\sigma,r}^F}(g)\right)\\
            &= \left(\mcF_{(\bg_{\sigma,r}^F)^*}\circ\mcF_{\bg_{\sigma,r}^F}(f)\right)*\left(\mcF_{(\bg_{\sigma,r}^F)^*}\circ\mcF_{\bg_{\sigma,r}^F}(g)\right).
        \end{align*}
         For any $\tilde f, \tilde g \in C\left((\bg_{\sigma,r}^F)^*\right)$, we can choose $f, g \in  C\left(\bg_{\sigma,r}^F\right)$ such that $\mcF_{\bg_{\sigma,r}^F}(f)=\tilde f$ and $\mcF_{\bg_{\sigma,r}^F}(g)=\tilde g$ (follows from $(i)$). Then, using the previous equation, we have $\mcF_{(\bg_{\sigma,r}^F)^*}(\tilde f \cdot \tilde g) = \mcF_{(\bg_{\sigma,r}^F)^*}(\tilde f) * \mcF_{(\bg_{\sigma,r}^F)^*}(\tilde g)$ and we are done. 
    \end{itemize}
    We have given the proofs for one of the results in each of $(ii)$, $(iii)$ and $(iv)$. The others follow from exactly similar calculations. Once we have these, the last statement follows immediately.  
\end{proof}

\end{document}
